\documentclass[11pt]{article}
\usepackage[margin=1in]{geometry}
\usepackage{hyperref}
\usepackage{enumitem}
\title{Minecraft Coding Fundamentals}
\author{Piter Garcia}
\date{\today}
\begin{document}
\maketitle

\section*{Grade and Class Match}
\textbf{Grade:} Grade 4\\
\textbf{Likely blocks:} Day 1 · 1:15-2:10 · Baris; Day 2 · 1:15-2:10 · Fowler; Day 3 · 1:15-2:10 · Schrank; Day 4 · 1:15-2:10 · Autore

\section*{Lesson Focus}
Minecraft Education setup, agent coding, challenge progression, and debugging.

\section*{Learning Targets}
\begin{itemize}[leftmargin=*]
\item Students enter the correct Minecraft coding world and follow setup steps.
\item Students use agent commands to solve a coding challenge.
\item Students debug a failed attempt and explain the fix.
\end{itemize}

\section*{Materials}
\begin{itemize}[leftmargin=*]
\item Student devices
\item Minecraft Education
\item Projected modeling screen
\end{itemize}

\section*{Suggested Lesson Flow}
\begin{enumerate}[leftmargin=*]
\item Open with the exact device, tool, or routine students need in order to start successfully.
\item Model one example task before students begin independent or partner work.
\item Give students time to build, code, design, or document their work while circulating for support.
\item Close with a short share-out, reflection, or debugging discussion.
\end{enumerate}

\section*{Source Notes}
This lesson packet is enriched with transcript-derived lesson notes, the school schedule roster, and official starting-point resources. See the paired Markdown guide for clickable links.

\bibliographystyle{plain}
\bibliography{minecraft-coding-fundamentals}
\end{document}
